\documentclass[exercice]{thedocument}%
\usepackage{texmaths-tikz}%
\startdatakeys
\source{originale}
\cycle{4}
\competencesDuSocle{CA, RE, CO}
\connaissancesRequises{A3a1,A3a4,A3a5}
\competencesTravaillees{A3d2}
\enddatakeys
\begin{document}%
Soit ABC un triangle rectangle en A. On sait que~:
\begin{itemize}
    \item AC = $2\times$AB
    \item Aire(ABC) = 50 cm$^2$
\end{itemize}
\begin{tikzpicture}[scale=2]
    \coordinate[label=below left:A](A)at(0,0);
    \coordinate[label=above left:B](B)at(0,1);
    \coordinate[label=below right:C](C)at(2,0);
    \coordinate[label=below:M](M)at(1,0);
    \coordinate(M1)at(1,0.05);
    \coordinate(M2)at(1,-0.05);
    \draw(A)--(B)--(C)--cycle;
    \draw(M1)--(M2);
    \tkzMarkSegments[mark=s||](A,B A,M M,C)
    \tkzMarkRightAngle[size=0.125](C,A,B)
\end{tikzpicture}\par\medskip\noindent
On rappelle la formule donnant l'aire d'un triangle:
$$\text{Aire}=\text{base}\times\text{hauteur}\div2$$
\begin{enumerate}
    \item En notant $x$ la longueur du segment [AB], exprimer l'aire du triangle ABC
    en fonction de $x$.
    \item Déterminer la longueur du segment [AB].
\end{enumerate}
\end{document}
